\chapter{Mathematical Modelling}

% -------------------------------------------------------
\section{Definition}
% -------------------------------------------------------

\begin{definition}[Mathematical Model]
A \textbf{mathematical model} is a representation of a real-life system or problem using mathematical equations and expressions, constructed in order to analyse, predict, and optimise the behaviour of that system.
\end{definition}

% -------------------------------------------------------
\section{Importance of Mathematical Modelling}
% -------------------------------------------------------

Mathematical modelling is central to engineering analysis and design. Its key roles include:

\begin{itemize}
  \item Simplifying complex systems into tractable mathematical forms.
  \item Predicting system behaviour under varying conditions without physical experiments.
  \item Aiding decision-making by providing quantitative insight.
  \item Optimising performance criteria such as cost, time, and efficiency.
  \item Enabling application across engineering, physics, economics, and the biological sciences.
\end{itemize}

% -------------------------------------------------------
\section{Types of Mathematical Models}
% -------------------------------------------------------

\subsection{Deterministic Model}

\begin{definition}
A \textbf{deterministic model} is one in which every set of variable states is uniquely determined by the model parameters and initial conditions.
\end{definition}

\begin{remark}
No randomness is involved; the same inputs always yield the same outputs.
\end{remark}

\subsection{Stochastic Model}

\begin{definition}
A \textbf{stochastic model} incorporates randomness or probabilistic elements, so the output is not uniquely fixed by the inputs.
\end{definition}

\begin{remark}
Repeated simulations with identical inputs may produce different outcomes.
\end{remark}

\subsection{Static Model}

\begin{definition}
A \textbf{static model} does not involve time as a variable. It describes the system at a single instant.
\end{definition}

\subsection{Dynamic Model}

\begin{definition}
A \textbf{dynamic model} describes how a system evolves over time.
\end{definition}

\begin{example}
The equation
\[
  \frac{dy}{dt} = 2y
\]
is a dynamic model of exponential growth. The unknown $y$ depends on time $t$.
\end{example}

\subsection{Linear Model}

\begin{definition}
A model is \textbf{linear} if all variables appear to the first power and are not multiplied together.
\end{definition}

\subsection{Non-Linear Model}

\begin{definition}
A model is \textbf{non-linear} if variables appear in higher powers, under transcendental functions, or as products of each other.
\end{definition}
